\documentclass[11pt,a4paper]{article}
\usepackage[T2A]{fontenc}
\usepackage[utf8]{inputenc}
\usepackage[english,russian]{babel}
\sloppy
\setlength{\emergencystretch}{2em}
\usepackage[margin=2.1cm]{geometry}
\usepackage{xcolor}
\usepackage{graphicx}
\usepackage{listings}
\usepackage{enumitem}
\usepackage{booktabs}
\usepackage{array}
\newcolumntype{L}[1]{>{\raggedright\arraybackslash}p{#1}}
\usepackage{amsmath}
\usepackage{tikz}
\usetikzlibrary{arrows.meta,positioning,calc,fit,shapes.geometric}
\usepackage{titlesec}
\usepackage{fancyhdr}
\usepackage[hidelinks]{hyperref}
\usepackage{caption}
\captionsetup{font=small,labelfont=bf}
% babel нь \begin{document}-д эдгээрийг дардаг тул AtBeginDocument ашиглана
\AtBeginDocument{%
  \renewcommand{\contentsname}{Гарчиг}%
  \renewcommand{\tablename}{Хүснэгт}%
  \renewcommand{\figurename}{Зураг}%
  \renewcommand{\refname}{Ашигласан материал}%
}

\definecolor{codebg}{RGB}{246,247,249}
\definecolor{accent}{RGB}{31,78,121}

\lstset{
  basicstyle=\ttfamily\footnotesize,
  backgroundcolor=\color{codebg},
  frame=single, rulecolor=\color{gray!40},
  keywordstyle=\color{accent}\bfseries,
  commentstyle=\color{gray}\itshape,
  breaklines=true, columns=fullflexible, keepspaces=true,
  extendedchars=true, inputencoding=utf8,
  literate=
    {А}{{\CYRA}}1
    {Б}{{\CYRB}}1
    {В}{{\CYRV}}1
    {Г}{{\CYRG}}1
    {Д}{{\CYRD}}1
    {Е}{{\CYRE}}1
    {Ё}{{\CYRYO}}1
    {Ж}{{\CYRZH}}1
    {З}{{\CYRZ}}1
    {И}{{\CYRI}}1
    {Й}{{\CYRISHRT}}1
    {К}{{\CYRK}}1
    {Л}{{\CYRL}}1
    {М}{{\CYRM}}1
    {Н}{{\CYRN}}1
    {О}{{\CYRO}}1
    {Ө}{{\CYROTLD}}1
    {П}{{\CYRP}}1
    {Р}{{\CYRR}}1
    {С}{{\CYRS}}1
    {Т}{{\CYRT}}1
    {У}{{\CYRU}}1
    {Ү}{{\CYRY}}1
    {Ф}{{\CYRF}}1
    {Х}{{\CYRH}}1
    {Ц}{{\CYRC}}1
    {Ч}{{\CYRCH}}1
    {Ш}{{\CYRSH}}1
    {Щ}{{\CYRSHCH}}1
    {Ъ}{{\CYRHRDSN}}1
    {Ы}{{\CYRERY}}1
    {Ь}{{\CYRSFTSN}}1
    {Э}{{\CYREREV}}1
    {Ю}{{\CYRYU}}1
    {Я}{{\CYRYA}}1
    {а}{{\cyra}}1
    {б}{{\cyrb}}1
    {в}{{\cyrv}}1
    {г}{{\cyrg}}1
    {д}{{\cyrd}}1
    {е}{{\cyre}}1
    {ё}{{\cyryo}}1
    {ж}{{\cyrzh}}1
    {з}{{\cyrz}}1
    {и}{{\cyri}}1
    {й}{{\cyrishrt}}1
    {к}{{\cyrk}}1
    {л}{{\cyrl}}1
    {м}{{\cyrm}}1
    {н}{{\cyrn}}1
    {о}{{\cyro}}1
    {ө}{{\cyrotld}}1
    {п}{{\cyrp}}1
    {р}{{\cyrr}}1
    {с}{{\cyrs}}1
    {т}{{\cyrt}}1
    {у}{{\cyru}}1
    {ү}{{\cyry}}1
    {ф}{{\cyrf}}1
    {х}{{\cyrh}}1
    {ц}{{\cyrc}}1
    {ч}{{\cyrch}}1
    {ш}{{\cyrsh}}1
    {щ}{{\cyrshch}}1
    {ъ}{{\cyrhrdsn}}1
    {ы}{{\cyrery}}1
    {ь}{{\cyrsftsn}}1
    {э}{{\cyrerev}}1
    {ю}{{\cyryu}}1
    {я}{{\cyrya}}1,
  xleftmargin=0.8em, framexleftmargin=0.8em,
  language=Java, showstringspaces=false
}
\titleformat{\section}{\Large\bfseries\color{accent}}{\thesection}{0.6em}{}
\titleformat{\subsection}{\large\bfseries}{\thesubsection}{0.6em}{}
\pagestyle{fancy}
\fancyhf{}
\fancyfoot[C]{\thepage}
\renewcommand{\headrulewidth}{0.4pt}
\setlength{\headheight}{14pt}
\setlength{\parskip}{3pt}
\fancyhead[L]{\small IV-2-7 --- Fork-Join Sorting Application}
\fancyhead[R]{\small GitHub repository}

\begin{document}
\begin{titlepage}
\centering
\vspace*{2.5cm}
{\large F.CSM302 --- Үйлдлийн систем · Бие даалт 1\par}
\vspace{1.4cm}
{\Huge\bfseries\color{accent} GitHub repository\\гарын авлага\par}
\vspace{0.5cm}
{\Large Fork-Join Sorting Application\par}
\vspace{0.3cm}
{\large Хуваарилалтын код: \textbf{IV-2-7} · Сэдэв №7\par}
\vspace{0.6cm}
{\large Кодыг судлах, шүүмжлэхэд зориулсан\par}
\vfill
{\large 2026 оны 10-р сар}
\vspace{2cm}
\end{titlepage}

\tableofcontents
\newpage

\section{Repository-гийн холбоос ба хувилбар}

\begin{center}
\small
\begin{tabular}{lL{10.0cm}}
\toprule
Талбар & Утга \\
\midrule
Холбоос & \texttt{https://github.com/<ХЭРЭГЛЭГЧ>/IV-2-7-ForkJoinSort} \\
Салбар (branch) & \texttt{main} \\
Хувилбар (tag) & \texttt{v1.0} --- 9-р долоо хоногт нийтэлсэн хувилбар \\
Лиценз / нууц өгөгдөл & Нууц түлхүүр, нууц үг, хувийн өгөгдөл \textbf{агуулаагүй} \\
Гадны сан & \textbf{Байхгүй}. Зөвхөн стандарт JDK. \\
\bottomrule
\end{tabular}
\end{center}

\begin{lstlisting}[language=bash]
git clone https://github.com/<ХЭРЭГЛЭГЧ>/IV-2-7-ForkJoinSort.git
cd IV-2-7-ForkJoinSort
git checkout v1.0          # үнэлгээнд ашиглах тогтсон хувилбар
\end{lstlisting}

\section{Фолдерын бүтэц}

\begin{lstlisting}[language=bash]
IV-2-7-ForkJoinSort/
+-- README.md              төслийн нэр, код, зорилго, Quick Start
+-- SUBMIT.md              илгээх заавар (багийн дотоод)
+-- build.sh / build.bat   компайл
+-- run.sh   / run.bat     ажиллуулах
+-- test.sh                тестийн багц
+-- src/forkjoinsort/      ЭХ КОД
+-- tests/                 58 тест
+-- examples/              оролтын жишээ файлууд
`-- docs/                  PDF-ууд, зураг, хэмжилтийн үр дүн
\end{lstlisting}

\begin{center}
\small
\begin{tabular}{lL{10.3cm}}
\toprule
Хэсэг & Даалгаврын шаардлага \\
\midrule
\texttt{README.md} & Төслийн нэр, код, зорилго, Quick Start \textbf{тийм} \\
\texttt{src/} & Эх код \textbf{тийм} \\
\texttt{tests/}, \texttt{examples/} & Үндсэн ба нэмэлт тестийн өгөгдөл \textbf{тийм} \\
\texttt{docs/} & Manual болон шаардлагатай зураг \textbf{тийм} \\
\texttt{v1.0} tag & 9-р долоо хоногт нийтэлсэн хувилбар \textbf{тийм} \\
\bottomrule
\end{tabular}
\end{center}

\section{Эх кодын файл бүрийн үүрэг}

\begin{center}
\small
\begin{tabular}{lrL{8.0cm}}
\toprule
Файл & Мөр & Үүрэг \\
\midrule
\multicolumn{3}{l}{\textbf{Алгоритмын цөм}} \\
\texttt{MergesortAction.java} & 45 & \texttt{RecursiveAction} --- тэнцүү хуваалт + merge \\
\texttt{QuicksortAction.java} & 65 & \texttt{RecursiveAction} --- 3 талт хуваалт, merge-гүй \\
\texttt{SequentialSort.java} & 140 & insertion sort, merge, \texttt{partition3},
pivot сонголт, нэг урсгалт жишиг \\
\midrule
\multicolumn{3}{l}{\textbf{Удирдлага ба хэмжилт}} \\
\texttt{SortRunner.java} & 75 & тохиргоо $\to$ үр дүн. \textbf{UI-ээс бүрэн хамааралгүй} \\
\texttt{RunConfig.java} & 55 & бүх параметр + оролтын шалгалт \\
\texttt{RunResult.java} & 30 & хэмжилт, мод, зөв эсэх \\
\texttt{DataGenerator.java} & 60 & 5 хэв маягийн өгөгдөл, текст оролт задлах \\
\texttt{Bench.java} & 80 & гурван хэмжилтийн туршилт (CSV) \\
\texttt{Main.java} & 120 & нэвтрэх цэг: GUI, \texttt{--headless}, \texttt{--bench}, \texttt{--render} \\
\midrule
\multicolumn{3}{l}{\textbf{Модны бүртгэл}} \\
\texttt{TraceRecorder.java} & 70 & \textbf{thread-safe} бүртгэл
(\texttt{AtomicInteger}, \texttt{ConcurrentLinkedQueue}) \\
\texttt{TraceNode.java} & 35 & модны нэг зангилаа = нэг \texttt{compute()} дуудлага \\
\midrule
\multicolumn{3}{l}{\textbf{Интерфейс}} \\
\texttt{ui/TreePanel.java} & 250 & \textbf{РЕКУРСИЙН МОД} зурах самбар \\
\texttt{ui/MainWindow.java} & 300 & цонх, удирдлага, статистик, legend \\
\bottomrule
\end{tabular}
\end{center}

\section{Кодыг ямар дарааллаар унших вэ}

Шүүмжлэгчид зориулсан санал болгох маршрут --- \textbf{15--20 минут}:

\begin{enumerate}[leftmargin=*,itemsep=3pt]
  \item \textbf{\texttt{MergesortAction.java}} (45 мөр) --- бүх зүйлийн гол.
        \texttt{compute()} функц нь: threshold шалгах $\to$ хуваах $\to$
        \texttt{invokeAll} $\to$ \texttt{merge}. Сурах бичгийн Project 3-ын
        шаардлага энд шууд харагдана.
  \item \textbf{\texttt{QuicksortAction.java}} (65 мөр) --- ижил бүтэц, гэхдээ
        ажил нь хуваалтын \emph{өмнө}. Гурван талт хуваалтын ач тус
        (\texttt{hasL}, \texttt{hasR} шалгалт) энд.
  \item \textbf{\texttt{SequentialSort.partition3()}} --- Dijkstra-гийн гурван
        талт гүйлт. Яагаад pivot-ыг урьдчилан зөөхгүй байгаа тайлбарыг
        коммент дотор уншина уу.
  \item \textbf{\texttt{SortRunner.run()}} --- туршилтын урсгал:
        жишиг хэмжилт $\to$ fork-join $\to$ \texttt{Arrays.sort}-той тулгах.
  \item \textbf{\texttt{TraceRecorder.java}} --- олон урсгалаас зэрэг бичигдэх
        бүртгэлийг цоожгүйгээр хэрхэн аюулгүй болгосон (6--7-р бүлгийн хэрэглээ).
  \item \textbf{\texttt{ui/TreePanel.layoutTree()}} --- зангилааны дугаар нь
        модны бүтцийг илэрхийлэхгүй (урсгалууд зэрэг үүсгэдэг тул) учраас
        эцгийн дугаараар дахин барьж DFS-ээр байрлуулдаг.
\end{enumerate}

\begin{center}
\small
\begin{tabular}{L{5.0cm}L{9.5cm}}
\toprule
Хайж байгаа зүйл & Хаанаас харах \\
\midrule
``\texttt{RecursiveAction} мөн үү?'' & \texttt{MergesortAction.java:14},
\texttt{QuicksortAction.java:16} \\
``\texttt{Comparable} мөн үү?'' & Хоёр файлын класс зарлалт:
\texttt{<T extends Comparable<? super T>>} \\
``Threshold хаана шалгагддаг вэ?'' & \texttt{compute()}-ийн эхний \texttt{if} \\
``Insertion sort мөн үү?'' & \texttt{SequentialSort.insertionSort()} \\
``Оролтын шалгалт хаана вэ?'' & \texttt{RunConfig.validate()} \\
``Зөв эсэхийг яаж шалгасан бэ?'' & \texttt{SortRunner.run()} доторх
\texttt{Arrays.equals(arr, expect)} \\
\bottomrule
\end{tabular}
\end{center}

\section{Нэг минутад шалгах}

\begin{lstlisting}[language=bash]
./build.sh && ./test.sh
# Хүлээгдэх: === ДҮН: 58 давсан, 0 унасан ===   (exit code 0)

./run.sh --headless -i "7,12,19,3,18,4,2,6,15,8" -t 2
# Хүлээгдэх: Зөв эсэх : ЗӨВ,  [2, 3, 4, 6, 7, 8, 12, 15, 18, 19]

./run.sh --render tree.png -n 400 -t 25 -a quick
# Дэлгэцгүй ч рекурсийн модыг зураг болгож харж болно
\end{lstlisting}

\section{Commit-ийн түүх ба багийн оролцоо}

Багийн гишүүн бүр өөрийн хувь нэмрийг commit-ээр баримтжуулсан.
Түүхийг харах:

\begin{lstlisting}[language=bash]
git log --oneline --graph --all
git shortlog -sn            # гишүүн бүрийн commit-ийн тоо
git log --stat v1.0         # хувилбарт орсон өөрчлөлтүүд
\end{lstlisting}

\begin{center}
\small
\begin{tabular}{llL{7.6cm}}
\toprule
Оюутны код & Нэр & Хариуцсан ажил \\
\midrule
B200910032 & Д.Дагиймаа & UI: \texttt{TreePanel} --- рекурсийн модны зураглал, зангилаа сонгох, алхамчилсан горим; \textbf{хэрэглэгчийн гарын авлага} \\
B221940034 & Д.Оюунжаргал & \texttt{TraceRecorder}, \texttt{TraceNode}, \texttt{DataGenerator} --- модны thread-safe бүртгэл, өгөгдөл үүсгэх; \textbf{README.md} \\
B241940014 & Г.Мягмарсүрэн & \texttt{TestSuite} --- 58 тест, \texttt{examples/}, хүчингүй оролтын боловсруулалт; \textbf{operational гарын авлага} \\
B241940127 & О.Энх & Багийн ахлагч. Fork/Join цөм: \texttt{MergesortAction}, \texttt{QuicksortAction}, \texttt{SequentialSort}, \texttt{SortRunner}; нэгтгэл, тайлан, \textbf{GitHub repo / tag / release} \\
B242270084 & Э.Бүрэнбаяр & \texttt{Bench.java} --- threshold / урсгал / хэв маягийн хэмжилт, \texttt{bench.txt}, үр дүнгийн дүгнэлт; \textbf{постын хавсралт} \\
\bottomrule
\end{tabular}
\end{center}

\section{Known issues}

\begin{enumerate}[nosep,leftmargin=*]
  \item Модны бүртгэл нь хурдад нөлөөлдөг тул хэмжилтийг
        \texttt{--no-trace} горимоор авна (\texttt{--bench} үүнийг автоматаар хийнэ).
  \item $n = 10^6$, threshold $= 10$ үед 200\,000 орчим зангилаа үүсч
        дүрслэл удаашрана. Дүрслэлд $n \le 5000$ хооронд ажиллуулахыг зөвлөнө.
  \item Зангилааны утгуудыг зөвхөн 48-аас бага мужид хадгална (санах ой).
  \item Хурдасгалт машины цөмийн тооноос хамаарна --- хөгжүүлэлтийн орчин
        2 цөмтэй байсан.
  \item \texttt{Comparable} шаардлагаас болж \texttt{int[]} биш
        \texttt{Integer[]} ашигласан (boxing-ийн зардалтай).
\end{enumerate}

\end{document}
